\documentclass[11pt]{article}
\usepackage[margin=2.1cm]{geometry}
\usepackage{booktabs}
\usepackage{amsmath}
\usepackage{enumitem}
\usepackage[T1]{fontenc}
\usepackage{lmodern}

\setlength{\parindent}{0pt}
\setlength{\parskip}{5pt}
\newcommand{\code}[1]{\texttt{#1}}
\newcommand{\answerbox}[1]{\vspace{2pt}\fbox{\parbox[c][#1][c]{\dimexpr\linewidth-2\fboxsep-2\fboxrule}{\ }}\vspace{4pt}}

\begin{document}

\begin{center}
{\Large\bfseries EE\,362 --- Pre-Lab 3}\\[2pt]
{\large Multiplexers}\\[6pt]
\begin{tabular}{@{}p{6.5cm}p{4cm}p{4cm}@{}}
Name: \hrulefill & Section: Mon / Tue & Date: \hrulefill
\end{tabular}
\end{center}

\vspace{2pt}\hrule\vspace{4pt}

{\small\textbf{Due at the start of the lab session.} Five points, about twenty
minutes, on paper.}

\vspace{4pt}\hrule

\section*{Question 1 --- Paper first}

A 32-bit 3-to-1 multiplexer has 98 inputs.

\textbf{(a)} How many rows would an exhaustive (Type E) truth table need? Write
the expression, you do not need to evaluate it.

\answerbox{1.2cm}

\textbf{(b)} Write the complete Type C table for it. Four rows. Remember that
\code{Sel = "11"} needs a \emph{defined} answer, not a blank.

\begin{center}
\begin{tabular}{clc@{\hskip 8pt}||@{\hskip 8pt}l@{\hskip 20pt}l}
\toprule
\# & \code{Sel1} & \code{Sel0} & $Y$ & note \\
\midrule
0 & 0 & 0 & \rule{2.2cm}{0.4pt} & \rule{4.5cm}{0.4pt} \\[4pt]
1 & 0 & 1 & \rule{2.2cm}{0.4pt} & \rule{4.5cm}{0.4pt} \\[4pt]
2 & 1 & 0 & \rule{2.2cm}{0.4pt} & \rule{4.5cm}{0.4pt} \\[4pt]
3 & 1 & 1 & \rule{2.2cm}{0.4pt} & \rule{4.5cm}{0.4pt} \\
\bottomrule
\end{tabular}
\end{center}

\textbf{(c)} Write the canonical equation for $Y_i$, then factor $\overline{S_1}$
out of the first two terms. Name the law you used.

\answerbox{2.6cm}

\textbf{One sentence:} look at your factored form. What familiar component is the
bracketed sub-expression?

\answerbox{1.2cm}

\section*{Question 2 --- Interface}

Write the VHDL \code{entity} for \code{Mux3}, including its \code{generic}.
Interface only.

\answerbox{4.4cm}

\section*{Question 3 --- Test plan}

You cannot test $2^{98}$ input combinations. State the coverage argument you will
use instead, then list the vectors.

\begin{center}
\begin{tabular}{@{}c@{\hskip 8pt}p{4.4cm}@{\hskip 8pt}p{8.2cm}@{}}
\toprule
\# & \code{Sel}, and the data on A / B / C & What this vector proves \\
\midrule
1 & \rule{4.2cm}{0.4pt} & \rule{8cm}{0.4pt} \\[8pt]
2 & \rule{4.2cm}{0.4pt} & \rule{8cm}{0.4pt} \\[8pt]
3 & \rule{4.2cm}{0.4pt} & \rule{8cm}{0.4pt} \\[8pt]
4 & \rule{4.2cm}{0.4pt} & \rule{8cm}{0.4pt} \\[8pt]
5 & \rule{4.2cm}{0.4pt} & \rule{8cm}{0.4pt} \\[8pt]
6 & \rule{4.2cm}{0.4pt} & \rule{8cm}{0.4pt} \\
\bottomrule
\end{tabular}
\end{center}

\textbf{Then answer both, in one or two sentences each:}

\textbf{(a)} Suppose someone hands you a ``multiplexer'' that ignores its selector
completely and always outputs \code{A}. Which of your vectors catches it? If none
do, fix your table before you come to lab.

\answerbox{1.8cm}

\textbf{(b)} Suppose someone forgets row 3, so the output latches. Which of your
vectors catches \emph{that}?

\answerbox{1.8cm}

\vfill
{\small\emph{Question 3 is the one that matters.} It is the same question as
Section 5 of your lab report.}

\end{document}
